\section{Restrictions on the number of players}
\label{sec:player-counts}

The distinction between an arbitrary typed subgroup and the complete
symmetry group of a payoff table is decisive here. We first give the
algebraic reductions, explicitly identifying the step that uses payoffs.
All games in this section have three strategies per player.

\begin{theorem}\label{thm:ternary-player-counts}
Let $G$ be the complete exact automorphism group of an $n$-player game,
with player projection $q(G)=S_n$.
\begin{enumerate}
\item If $n\geq3$, $3\nmid n$, and the game is standard symmetric, then
$q:G\to S_n$ splits.
\item If $n\geq5$ and $n\neq6$, then $q$ splits, without any standardness
assumption.
\end{enumerate}
Consequently, in conjunction with the three-player splitting result and the
six-player construction, six is the only possible player count for a
standard symmetric ternary game with full player image and nonsplitting
projection.
\end{theorem}

The conclusion about the unique possible player count uses the classical
Schur multiplier theorem for alternating groups. That group-theoretic
input is established literature, rather than a new assertion about games.

Identify each strategy set with $\mathbb F_3$. Every local permutation is
uniquely affine, $x\mapsto ax+b$ with $a\in\{1,-1\}$. Thus the ambient
typed group is
\[
 W=R\rtimes V,\qquad R=\mathbb F_3^n,\qquad
 V=\mathbb F_2^n\rtimes S_n.
\]
Let $\rho:W\to V$ forget translations, and put
$B=\rho(G)$, $K=B\cap\mathbb F_2^n$, and $N=G\cap R$.
Write $T_\ell$ and $E_\ell$ for the constant line and the coordinate-sum-zero
subspace of $\mathbb F_\ell^n$, respectively.

\begin{lemma}\label{lem:ternary-sign-kernels}
Suppose $n\geq3$ and $q(G)=S_n$.
\begin{enumerate}
\item $K$ is one of $0,T_2,E_2,\mathbb F_2^n$.
\item If $K\supseteq E_2$ and $G$ is a complete game automorphism group,
the game is fully symmetric under a matching.
\item If $K=T_2$, a local translation conjugates $G$ to $N\rtimes B$.
Its player projection splits for $n\geq5$, and also for $n=3,4$ when
$G$ contains a transitive strategy-trivial subgroup.
\end{enumerate}
\end{lemma}
\begin{proof}
Conjugation makes $K$ a permutation-invariant subspace. If it contains a
nonconstant vector, subtracting its image under a transposition of unequal
entries gives a coordinate-pair vector. Its conjugates span $E_2$.
Otherwise $K\subseteq T_2$. This proves the first assertion.

Suppose $K\supseteq E_2$. The coordinate characters of the even-sign
group $E_2$ are distinct and nontrivial for $n\geq3$. Character averaging,
valid because $|E_2|$ is invertible in $\mathbb F_3$, projects any vector
of the invariant subspace $N$ onto each coordinate line. Player
transitivity therefore gives $N=0$ or $N=R$. If $N=R$, translation
invariance makes every payoff constant; full player image equates these
constants, giving the full ambient group.

If $N=0$, $G$ is the graph of a cocycle $c:B\to R$. More generally,
for any such cocycle and a normal subgroup $L$ of order prime to three,
put $w=|L|^{-1}\sum_{l\in L}c(l)$. The cocycle identity gives
$c(l)=w-lw$ on $L$. Translation conjugation makes $c$ vanish on $L$;
normality then implies $c(b)$ is fixed by $L$ for every $b\in B$.
Apply this to $L=E_2$. Its fixed space is zero, so the conjugated group
is exactly $B$, acting by signed coordinate permutations.

If $K=\mathbb F_2^n$, cancellation of signs gives every bare player
permutation. If $K=E_2$, the quotient group is the graph of a homomorphism
$S_n\to\mathbb F_2$, either zero or parity. The zero case again contains
all bare permutations. In the parity case, on a profile containing zero,
any sign change has the same effect as an even sign change: correct its
parity at a zero coordinate. On all-nonzero profiles, $E_2$ has two orbits,
determined by the product of the coordinates. For a focal player $i$,
a transposition fixing $i$ has an odd-sign lift, which interchanges these
two orbits and equates their values of $u_i$. Such a transposition exists
because $n\geq3$. Thus all independent sign changes preserve every payoff.
Cancelling signs in the given lifts now supplies all bare permutations.
Undoing conjugation proves full symmetry under a matching. This is the
step requiring actual payoffs; it is false for arbitrary subgroups.

For $K=T_2$, the image of $G$ in $(R/N)\rtimes B$ is a cocycle graph.
The central order-two subgroup $T_2$ acts as $-1$ on $R/N$. The same
averaging argument kills this cocycle, since the fixed space of $-1$ is
zero in characteristic three. Lifting the conjugating translation and
cancelling elements of $N$ gives $G=N\rtimes B$.

It remains to split $B\to S_n$. A lift of a transposition squares to a
vector supported on its two moved coordinates and belonging to $T_2$;
for $n\geq3$ this vector is zero. Choose involutory lifts of the adjacent
transpositions. For $n\geq5$, commutators of disjoint lifts are supported
on at most four coordinates and constant, so vanish. The cubes of adjacent
products lie in $T_2$; multiplying successive lifts by the central constant
flip kills these braid defects in order along the adjacent chain. These
are the Coxeter relations of $S_n$, giving a section. For $n=3$ there
are no distant relations, so the same argument applies.

For $n=4$, the distant commutator is either zero for every disjoint pair,
or the constant flip for every such pair. The alternative is independent
of lift choices and preserved by conjugation. In the former case the
Coxeter argument applies. In the latter, every lift of a double
transposition has order four. A transitive subgroup of $S_4$ contains a
double transposition: a Sylow two-subgroup has order at least four and
the subgroups of that size in $S_4$ have such an element. A transitive
strategy-trivial subgroup of $G$ maps to a strategy-trivial subgroup of
$B$ with the same player image; its projection is injective. It would
therefore supply an involutory lift of that double transposition, a
contradiction. Standardness excludes the latter alternative.
\end{proof}

\begin{lemma}\label{lem:ternary-residual-extension}
For $n\geq5$, if $G$ is a complete game group with $q(G)=S_n$ and
nonsplitting projection, then, after local relabeling,
\[
 K=0,\qquad N=T_3,\qquad |G|=3n!.
\]
Even player permutations centralize $N$, and odd ones invert it.
Furthermore $3\mid n$.
\end{lemma}
\begin{proof}
The preceding lemma excludes nonzero $K$. Thus $B\cong S_n$. A
point-stabilizer acts on its local sign fibre by a character
$S_{n-1}\to C_2$, either trivial or parity. Coherent local sign
relabeling therefore makes the action of $B$ on $R$ either the ordinary
permutation action or its common-parity twist. To see this explicitly,
choose player transports $\tau_i$ from a reference player to $i$.
Relabeling by their local components expresses the component of $p$
through the stabilizer element $\tau_{p(i)}^{-1}p\tau_i$. For the parity
character, a further relabeling by the signs of $\tau_i$ leaves the common
factor $(-1)^{\epsilon(p)}$.

The invariant subspaces of $R$ are $0,T_3,E_3,R$: a nonconstant vector,
subtracted from its transposed image, supplies a nonzero multiple of
$e_i-e_j$, whose conjugates span $E_3$. A parity twist changes no invariant
subspaces. If $3\nmid n$, the equivariant projection
\[
 v\longmapsto n^{-1}\Bigl(\sum_i v_i\Bigr)(1,\ldots,1)
\]
gives $R=T_3\oplus E_3$. Each possible $N$ then has an invariant
complement. The graph cocycle $B\to R/N$ lifts through that complement
to a cocycle $B\to R$, whose graph lies in $G$ and supplies a section.
This proves the divisibility assertion and also the first conclusion of
Theorem~\ref{thm:ternary-player-counts}, using the preceding lemma at $n=4$.

For the remaining cases, $N=0$ immediately splits the projection, while
$N=R$ makes the actual group the full ambient group. If $N=E_3$, the
quotient $R/N$ is one-dimensional with trivial or sign action. Its cocycle
vanishes on perfect $A_n$, and descends to $C_2$. For trivial action the
involution relation makes this cocycle zero; for sign action every value
is a coboundary, since two is invertible in $\mathbb F_3$. Thus translation
conjugation gives $E_3\rtimes B$, again split.

Suppose $N=T_3$ and the action is ordinary. Any transposition lift
$(v,p)$ has square $cz$, where $z=(1,\ldots,1)$. At every fixed
coordinate, $2v_i=c$. Subtracting $(c/2)z$ from its translation makes it
an involution with zero translation at all fixed coordinates. Choose such
lifts for adjacent transpositions. Distant commutators have translation
supported on at most four players; adjacent-product cubes have translation
supported on at most three. Both belong to $T_3$, so vanish for $n\geq5$.
The Coxeter relations give a section. Only the parity-twisted action with
$N=T_3$ remains, proving all assertions.
\end{proof}

\begin{lemma}\label{lem:alternating-multiplier-splitting}
For $n\geq8$, an extension of $S_n$ by $C_3$ in which even permutations
centralize the kernel and odd permutations invert it has a section.
\end{lemma}
\begin{proof}
The classical Schur multiplier theorem gives $|M(A_n)|=2$ for $n\geq8$;
see Lassueur and Mazza~\cite[Section~3]{LassueurMazza2015}, which recalls the theorem and its
standard references. Set $H=q^{-1}(A_n)$. The universal central extension
of the perfect group $A_n$ maps to this central $C_3$ extension. Its
central kernel, of order two, has trivial image in $C_3$. Its image
therefore supplies a perfect complement $D\cong A_n$, so
$H=C_3\times D$ and $D=H'$. Thus $D$ is normal in $G$.

If $t$ lifts a transposition, then $t^2\in C_3$. Conjugation by $t$
both fixes its own square and inverts $C_3$, so $t^2=(t^2)^{-1}$,
forcing $t^2=1$. The subgroup $\langle D,t\rangle$ now has order $n!$,
intersects the kernel trivially, and projects onto $S_n$. It is the
required section.
\end{proof}

Combining Lemmas~\ref{lem:ternary-residual-extension} and
\ref{lem:alternating-multiplier-splitting} excludes every $n\geq8$;
$n=5,7$ are excluded by $3\nmid n$. This proves the second assertion of
Theorem~\ref{thm:ternary-player-counts}. For standard games the one-player
case is trivial, the two-player standard subgroup itself supplies a
section, the three-player result was proved earlier, and the first
assertion covers four players. The explicit six-player game attains the
remaining possibility.

Actualness cannot be omitted. For even $n\geq4$ the ternary typed subgroup
\[
 D_n=\{(b,p)\in\mathbb F_2^n\rtimes S_n:
              \textstyle\sum_i b_i=\epsilon(p)\},
 \qquad x_i\mapsto(-1)^{b_i}x_i,
\]
contains bare $A_n$ and is therefore standard and full-image. Every lift
of an $n$-cycle has $n$th power the nontrivial constant flip, hence order
$2n$. Its projection cannot split. Lemma~\ref{lem:ternary-sign-kernels}
explains why this group cannot be the complete automorphism group of an
invariant payoff table.
